% GENERATED FILE. Edit canonical lessons, then run npm run generate:books.
% canonical-source-hash: fnv1a64:1918147d245e4ac8
% canonical-lessons: ZH-C427-gebi, ZH-C427-laiwang, ZH-W427-wang-toward, ZH-C427-zizhucan, ZH-C427-chuangdan, ZH-C427-kaiwang

\chapter{Comings and Goings}
\label{ch:zh-comings-and-goings}
\hlchaptermodality{427}

\begin{chapteropening}
\textbf{By the end of this chapter:} \emph{I can talk about the room next door, traffic coming and going, the buffet, the sheets and a bus's route.}

Complete the last lesson of chapter 427: I can talk about the room next door, traffic coming and going, the buffet, the sheets and a bus's route.
\end{chapteropening}

\section[gébì]{gébì — next door}

\label{lesson:ZH-C427-gebi}



\begin{quote}
Before the new one: say the Mandarin for the hotel star rating, then the Mandarin for a power socket.
\end{quote}

\subsection*{You'll want to know: gébì}
\emph{gébì} — \textquotedblleft{}next door\textquotedblright{}.

\begin{quote}
\emph{wǒ zhù zài gébì} — I'm staying next door.
\end{quote}

Its characters come later; for now it is a word for the ear.

\begin{cousinweb}
\emph{gébì} is built from \emph{gé} 'to separate' and \emph{bì} 'a wall', both inherited from Middle Chinese, \emph{keak} and \emph{pek} (Guangyun, in Baxter's transcription), which glosses the first 'to block' and the second 'a wall'.
\end{cousinweb}

\begin{grammarlens}[title={What you've built}]
One more word for getting around.
\end{grammarlens}

\subsection*{Your turn}
\begin{itemize}
\raggedright
  \item \emph{Say it:} \emph{gébì}
  \item \emph{Say it:} \emph{gébì}, once more
  \item \emph{Say it:} \emph{wǒ zhù zài gébì} — I'm staying next door.
  \item \emph{Recall:} say the Mandarin for the hotel star rating, then the Mandarin for a power socket, then say \emph{gébì} again
\end{itemize}

\begin{tcolorbox}[breakable,colback=teal!4,colframe=teal!35!black,title={Before you move on}]
What does gébì mean? (\textquotedblleft{}Next door\textquotedblright{}.) Say it once more.
\end{tcolorbox}
\section[\texorpdfstring{\zh{来往} (láiwǎng)}{来往 (láiwǎng)}]{\zh{来往} (láiwǎng) — to come and go}

\label{lesson:ZH-C427-laiwang}



\begin{quote}
Before the new one: say the Mandarin for a power socket, then the Mandarin for right next door.
\end{quote}

\subsection*{You'll want to know: \zh{来往}}
\textbf{\zh{来往}} — \emph{láiwǎng} — \textquotedblleft{}to come and go\textquotedblright{}.

\begin{quote}
\textbf{\zh{路口来往的车很快你要小心}} \emph{lùkǒu láiwǎng de chē hěn kuài nǐ yào xiǎoxīn} — The cars coming and going at the junction are fast: be careful.
\end{quote}

The new character, in pieces: \zh{往} — eight strokes in two pieces: on the left the left piece of \zh{很}, three strokes, two short strokes falling to the left and a vertical; on the right a dot, two bars, a vertical through them and a long bar at the foot. The next lesson puts it in your hand.

\begin{cousinweb}
\textbf{\zh{来往}} (\emph{láiwǎng}) is built from \zh{来} 'to come' and \zh{往} 'to go', both inherited from Middle Chinese, \emph{loj} and \emph{hjwangX} (Guangyun, in Baxter's transcription), which glosses the second 'to go'.
\end{cousinweb}

\begin{grammarlens}[title={What you've built}]
One more word for getting around.
\end{grammarlens}

\subsection*{Your turn}
\begin{itemize}
\raggedright
  \item \emph{Say it:} \emph{láiwǎng}
  \item \emph{Say it:} \emph{láiwǎng}, once more
  \item \emph{Say it:} \emph{lùkǒu láiwǎng de chē hěn kuài nǐ yào xiǎoxīn} — The cars coming and going at the junction are fast: be careful.
  \item \emph{Recall:} say the Mandarin for a power socket, then the Mandarin for right next door, then say \emph{láiwǎng} again
\end{itemize}

\begin{tcolorbox}[breakable,colback=teal!4,colframe=teal!35!black,title={Before you move on}]
What does \zh{来往} mean? (\textquotedblleft{}To come and go\textquotedblright{}.) Say it once more.
\end{tcolorbox}
\section[\texorpdfstring{\zh{往} (wǎng)}{往 (wǎng)}]{\zh{往} — the left piece of \zh{很}, then a dot over three bars}

\label{lesson:ZH-W427-wang-toward}



\begin{quote}
The word first, then the character inside it.

\begin{itemize}
\raggedright
  \item \emph{Recall:} say \emph{láiwǎng} — \textbf{R1}, one lesson back, before it decays
\end{itemize}

\begin{quote}
\zh{很} \zh{来}
\end{quote}

\emph{Read it:} each one aloud, with nothing said first — \emph{hěn}, \emph{lái}
\end{quote}

\subsection*{Script — the left piece of \zh{很}, then a dot over three bars}
\hlblockfigure{\detokenize{figures/ZH-W427-wang-toward-filmstrip.pdf}}{How \zh{往} is written, stroke by stroke}

\begin{quote}
\zh{往}
\end{quote}

\textbf{Two pieces, left before right.} Count the pieces, not the strokes.

\textbf{The left piece, three strokes}, as in \zh{很}: two short strokes falling to the left, one under the other, then a vertical down from the second.

\textbf{The right piece, five strokes.} A dot on top. A bar, then a second bar. A vertical down through them. Then a long bar across the foot.

3 + 5 = eight.

\subsection*{Your turn}
\begin{itemize}
\raggedright
  \item \emph{Write it:} \zh{往}, all eight, with the model in front of you
  \item \emph{Cover:} write it again from memory, and count the pieces, not the strokes
\end{itemize}

\begin{tcolorbox}[breakable,colback=teal!4,colframe=teal!35!black,title={Before you move on}]
How many strokes in \zh{往}? (\textbf{Eight}.) What is the left piece? (\textbf{The left piece of \zh{很}}.) What closes the right piece at the foot? (\textbf{A long bar}.)

Source: \href{https://raw.githubusercontent.com/chanind/hanzi-writer-data/68d10a4b21150cae5e1ebbd223eed289cf32d90c/data/\%E5\%BE\%80.json}{Hanzi Writer Data \zh{往}, pinned stroke paths and medians}.
\end{tcolorbox}
\section[zìzhùcān]{zìzhùcān — buffet}

\label{lesson:ZH-C427-zizhucan}



\begin{quote}
Before the new one: say the Mandarin for right next door, then the Mandarin for to come and go.
\end{quote}

\subsection*{You'll want to know: zìzhùcān}
\emph{zìzhùcān} — \textquotedblleft{}buffet\textquotedblright{}.

\begin{quote}
\emph{zhè jiā jiǔdiàn de zǎofàn shì zìzhùcān} — Breakfast at this hotel is a buffet.
\end{quote}

Its characters come later; for now it is a word for the ear.

\begin{cousinweb}
\emph{zìzhùcān} is built from \emph{zì} 'oneself', \emph{zhù} 'to help' and \emph{cān} 'a meal', all inherited from Middle Chinese, \emph{dzijH}, \emph{dzrjoH} and \emph{tshan} (Guangyun, in Baxter's transcription), which glosses the third 'to swallow'.
\end{cousinweb}

\begin{grammarlens}[title={What you've built}]
One more word for getting around.
\end{grammarlens}

\subsection*{Your turn}
\begin{itemize}
\raggedright
  \item \emph{Say it:} \emph{zìzhùcān}
  \item \emph{Say it:} \emph{zìzhùcān}, once more
  \item \emph{Say it:} \emph{zhè jiā jiǔdiàn de zǎofàn shì zìzhùcān} — Breakfast at this hotel is a buffet.
  \item \emph{Recall:} say the Mandarin for right next door, then the Mandarin for to come and go, then say \emph{zìzhùcān} again
\end{itemize}

\begin{tcolorbox}[breakable,colback=teal!4,colframe=teal!35!black,title={Before you move on}]
What does zìzhùcān mean? (\textquotedblleft{}Buffet\textquotedblright{}.) Say it once more.
\end{tcolorbox}
\section[chuángdān]{chuángdān — bed sheet}

\label{lesson:ZH-C427-chuangdan}



\begin{quote}
Before the new one: say the Mandarin for to come and go, then the Mandarin for a buffet.
\end{quote}

\subsection*{You'll want to know: chuángdān}
\emph{chuángdān} — \textquotedblleft{}bed sheet\textquotedblright{}.

\begin{quote}
\emph{wǒ xiǎng huàn chuángdān} — I'd like the sheets changed.
\end{quote}

Its characters come later; for now it is a word for the ear.

\begin{cousinweb}
\emph{chuángdān} is built from \emph{chuáng} 'a bed' and \emph{dān} 'a single layer', both inherited from Middle Chinese, \emph{dzrjang} and \emph{tan} (Guangyun, in Baxter's transcription).
\end{cousinweb}

\begin{grammarlens}[title={What you've built}]
One more word for getting around.
\end{grammarlens}

\subsection*{Your turn}
\begin{itemize}
\raggedright
  \item \emph{Say it:} \emph{chuángdān}
  \item \emph{Say it:} \emph{chuángdān}, once more
  \item \emph{Say it:} \emph{wǒ xiǎng huàn chuángdān} — I'd like the sheets changed.
  \item \emph{Recall:} say the Mandarin for to come and go, then the Mandarin for a buffet, then say \emph{chuángdān} again
\end{itemize}

\begin{tcolorbox}[breakable,colback=teal!4,colframe=teal!35!black,title={Before you move on}]
What does chuángdān mean? (\textquotedblleft{}Bed sheet\textquotedblright{}.) Say it once more.
\end{tcolorbox}
\section[\texorpdfstring{\zh{开往} (kāiwǎng)}{开往 (kāiwǎng)}]{\zh{开往} (kāiwǎng) — to be bound for a place}

\label{lesson:ZH-C427-kaiwang}



\begin{quote}
Before the new one: say the Mandarin for a buffet, then the Mandarin for a bed sheet.
\end{quote}

\subsection*{You'll want to know: \zh{开往}}
\textbf{\zh{开往}} — \emph{kāiwǎng} — \textquotedblleft{}to be bound for a place\textquotedblright{}.

\begin{quote}
\textbf{\zh{开往机场的车在酒店门口}} \emph{kāiwǎng jīchǎng de chē zài jiǔdiàn ménkǒu} — The bus to the airport is at the hotel entrance.
\end{quote}

\begin{cousinweb}
\textbf{\zh{开往}} (\emph{kāiwǎng}) is built from \zh{开} 'to drive off' and \zh{往} 'to go towards', both inherited from Middle Chinese, \emph{khoj} and \emph{hjwangX} (Guangyun, in Baxter's transcription), which glosses the second 'to go'.
\end{cousinweb}

\begin{grammarlens}[title={What you've built}]
One more word for getting around.
\end{grammarlens}

\subsection*{Your turn}
\begin{itemize}
\raggedright
  \item \emph{Say it:} \emph{kāiwǎng}
  \item \emph{Say it:} \emph{kāiwǎng}, once more
  \item \emph{Say it:} \emph{kāiwǎng jīchǎng de chē zài jiǔdiàn ménkǒu} — The bus to the airport is at the hotel entrance.
  \item \emph{Recall:} say the Mandarin for a buffet, then the Mandarin for a bed sheet, then say \emph{kāiwǎng} again
\end{itemize}

\begin{tcolorbox}[breakable,colback=teal!4,colframe=teal!35!black,title={Before you move on}]
What does \zh{开往} mean? (\textquotedblleft{}To be bound for a place\textquotedblright{}.) Say it once more.
\end{tcolorbox}
